\documentclass[a0paper, portrait]{tikzposter}

% Загрузка пакетов
\usepackage[utf8]{inputenc}
\usepackage[T2A]{fontenc}
\usepackage[russian]{babel}
\usepackage{graphicx}
\usepackage{amsmath, amssymb}
\usepackage{tikz}
\usepackage{array}  % Для улучшенных таблиц

% Метаданные постера
\title{Пример постера с tikzposter}
\author{Викторов Егор}
\institute{РУДН}
\titlegraphic{logo.png}  % Логотип (если есть)

% Выбор цветовой схемы и темы
\usetheme{Board}
\usebackgroundstyle{Rays}
\usecolorstyle{BlueGray}

\begin{document}

% Заголовок постера
\maketitle

% Основная сетка: две колонки
\begin{columns}

  % Первая колонка (левая)
  \column{0.45}
  
  % Блок с заголовком
  \block{Введение}{
    Это пример блока с заголовком. Здесь можно разместить текст, формулы, изображения.
    
    Пример формулы:
    \[
      E = mc^2
    \]
    
    % Вставка изображения (в окружении center)
    \begin{center}
      \includegraphics[width=0.8\linewidth]{example-image-a}
    \end{center}
  }
  
  % Блок без заголовка (пустые braces)
  \block{}{
    Этот блок без заголовка. Первые фигурные скобки оставлены пустыми.
    
    \begin{itemize}
      \item Пункт списка 1
      \item Пункт списка 2
      \item Пункт списка 3
    \end{itemize}
  }

  % Вторая колонка (правая)
  \column{0.55}
  
  % Блок с заметкой
  \block{Раздел с заметкой}{
    Содержимое блока. Добавим вертикальный отступ для заметки:
    
    \vspace{4cm}  % Создаём место для заметки
    
    % Вставка таблицы
    \begin{center}
      \begin{tabular}{|c|c|c|}
        \hline
        Столбец 1 & Столбец 2 & Столбец 3 \\
        \hline
        1 & 2 & 3 \\
        4 & 5 & 6 \\
        \hline
      \end{tabular}
    \end{center}
  }
  
  % Заметка рядом с блоком
  \note[
    targetoffsetx=3cm,   % Смещение по горизонтали
    targetoffsety=-2cm, % Смещение по вертикали
    width=0.3\linewidth, % Ширина заметки
    color=red!20         % Цвет фона заметки
  ]{
    Это заметка, прикреплённая к блоку. Её положение задаётся параметрами \\
    \texttt{targetoffsetx} и \texttt{targetoffsety}.
  }

\end{columns}

% Дополнительный блок на всю ширину (без колонок)
\block{Полный блок}{
  Этот блок занимает всю ширину постера. Здесь можно разместить:
  * текст;
  * изображения (в \texttt{center});
  * таблицы;
  * библиографию (если нужно).

  \begin{center}
    \includegraphics[width=0.4\linewidth]{example-image-b}
  \end{center}
}

\end{document}
